% Look and feel for the Step 1 notes. The main text is in notes.tex.
% Only the typeface block changes between design variants.

\usepackage{fontspec}
\usepackage[a4paper, left=2.7cm, textwidth=13.2cm, marginparsep=0.7cm,
  marginparwidth=4.0cm, top=3.0cm, bottom=3.2cm]{geometry}
\usepackage{xcolor}
\usepackage{booktabs}
\usepackage{tabularx}
\usepackage{amsmath}
\usepackage{amssymb}
\usepackage{enumitem}
\usepackage{listings}
\usepackage{tcolorbox}
\tcbuselibrary{skins}
\usepackage{marginnote}
\usepackage{tikz}
\usetikzlibrary{arrows.meta, positioning, calc, fit, shapes.geometric}
\usepackage{caption}
\usepackage{float}
\usepackage[hidelinks]{hyperref}

% ---- typeface (the design axis) -----------------------------------
\setmainfont{EBGaramond}[Extension=.otf, UprightFont=*-Regular, BoldFont=*-Bold, ItalicFont=*-Italic, BoldItalicFont=*-BoldItalic]
\newcommand{\fontlabel}{EB Garamond}

\setsansfont{texgyreheros}[Extension=.otf, UprightFont=*-regular,
  BoldFont=*-bold, ItalicFont=*-italic, BoldItalicFont=*-bolditalic,
  Scale=0.88]
\setmonofont{lmmono10}[Extension=.otf, UprightFont=*-regular,
  Scale=0.92]

% ---- colours: calm neutrals, one accent ----------------------------
\definecolor{ink}{HTML}{2E2622}
\definecolor{muted}{HTML}{7A6D64}
\definecolor{line}{HTML}{DDD3C7}
\definecolor{paper}{HTML}{F6F1EA}
\definecolor{accent}{HTML}{B5583A}
\definecolor{accentbg}{HTML}{F5E6DF}

\color{ink}
\linespread{1.12}
\setlength{\parindent}{1.2em}
\setlength{\parskip}{0pt}
\emergencystretch=3em
\frenchspacing

% ---- headings ------------------------------------------------------
\makeatletter
\renewcommand\section{\@startsection{section}{1}{0pt}%
  {3.2ex plus 1ex minus .2ex}{1.6ex}%
  {\normalfont\Large\color{ink}}}
\renewcommand\subsection{\@startsection{subsection}{2}{0pt}%
  {2.6ex plus .8ex minus .2ex}{1.0ex}%
  {\normalfont\large\itshape\color{ink}}}
\renewcommand\@seccntformat[1]{%
  \llap{\color{accent}\csname the#1\endcsname\hspace{0.9em}}}
\makeatother

% ---- running head and foot ----------------------------------------
\usepackage{fancyhdr}
\pagestyle{fancy}
\fancyhf{}
\renewcommand{\headrulewidth}{0pt}
\fancyhead[L]{\footnotesize\scshape\color{muted}%
  \MakeLowercase{Queries, scores and filters}}
\fancyhead[R]{\footnotesize\color{muted}\thepage}
\fancyheadoffset[R]{0pt}

% ---- margin notes: a short explanation of a term ------------------
\newcommand{\term}[2]{\emph{#1}\marginnote{\footnotesize\raggedright
  \color{muted}\textbf{\color{ink}#1:} #2}}
\renewcommand{\marginfont}{\footnotesize}

% ---- boxes ---------------------------------------------------------
\newtcolorbox{keypoint}{colback=accentbg, colframe=accentbg,
  boxrule=0pt, arc=0pt, left=10pt, right=10pt, top=7pt, bottom=7pt,
  borderline west={2pt}{0pt}{accent}}
\newtcolorbox{example}[1]{colback=white, colframe=line, boxrule=0.5pt,
  arc=0pt, left=10pt, right=10pt, top=6pt, bottom=6pt,
  fonttitle=\small\scshape, coltitle=muted, colbacktitle=white,
  title={\MakeLowercase{#1}}, titlerule=0.4pt, titlerule style=line}

% ---- code ----------------------------------------------------------
\lstset{
  basicstyle=\ttfamily\small,
  backgroundcolor=\color{paper},
  frame=none,
  xleftmargin=0.7em, xrightmargin=0.7em,
  aboveskip=0.9em, belowskip=0.9em,
  columns=fullflexible, keepspaces=true,
  showstringspaces=false,
  keywordstyle=\color{accent},
  commentstyle=\color{muted}\itshape,
  literate={×}{{$\times$}}1 {→}{{$\rightarrow$}}1
           {æ}{{\ae}}1 {ø}{{\o}}1 {å}{{\aa}}1 {ü}{{\"u}}1
}
\newcommand{\code}[1]{\texttt{#1}}

% ---- captions and tables -------------------------------------------
\captionsetup{font={small}, labelfont={sc,color=muted},
  textfont={color=muted}, labelsep=quad, justification=raggedright,
  singlelinecheck=false}
\renewcommand{\arraystretch}{1.18}

% ---- figure styles -------------------------------------------------
\tikzset{
  box/.style={draw=line, fill=white, rounded corners=3pt,
    font=\small, align=center, inner sep=6pt, minimum height=8mm},
  hot/.style={box, draw=accent, fill=accentbg},
  flow/.style={-{Stealth[length=2.2mm]}, draw=muted, line width=0.6pt},
  hotflow/.style={-{Stealth[length=2.2mm]}, draw=accent,
    line width=0.8pt},
  lbl/.style={font=\footnotesize, text=muted, align=center},
}

\raggedbottom
